% Conclusions and recommendations (BGU \S 10.7.2.9).
\section{Conclusions and Recommendations}
\label{sec:conclusions}

This section summarizes the project, compares the achieved performance to
the preliminary specification, discusses the strengths and weaknesses of
the chosen approach, and lists the recommendations for any continuation.

\subsection{Work summary}
The project delivered a reproducible 402-cell training campaign across
three architectures (ResNet50, DenseNet121, TransNeXt-tiny), two datasets
(CIFAR-10, MNIST), five degradation levels, six experimental phases, and
a 48-replicate multi-seed audit. The campaign is anchored to a
deterministic 256-image PSNR / SSIM probe per pipeline, so the cross-phase
comparison sits on a physical scale. Every value reported in this report
is reproducible from a clean checkout using the scripts under
\texttt{Final\_Report/scripts/}.

\subsection{Plan vs achieved}
Table \ref{tab:plan_vs_actual} compares the preliminary specification
acceptance criteria against the measured outcome. Figure
\ref{fig:gantt} compares the planned schedule against the actual schedule.

\begin{table}[htbp]
\centering
\caption{Acceptance criteria from the preliminary specification vs the measured outcome.}
\label{tab:plan_vs_actual}
\setlength{\tabcolsep}{4pt}
\renewcommand{\arraystretch}{1.15}
\begin{tabular}{p{0.30\textwidth}p{0.30\textwidth}p{0.30\textwidth}}
\toprule
Criterion & Plan & Actual \\
\midrule
Top-1 accuracy on the moderate-degradation operating point
& $\geq 80\%$ on CIFAR-10 L3
& Met at L3 on MNIST (every backbone). Met at L1 and L2 on CIFAR-10. Not met at L3 on CIFAR-10 (best $\sim 74\%$); not met at L4 / L5 on either dataset. \\
\midrule
Accuracy drop ceiling
& $\leq 15$ \Mpp{} drop relative to clean
& Met up to L2 on CIFAR-10 and up to L3 on MNIST. Not met at L4 / L5 on either dataset. \\
\midrule
Architectural winner
& A specific model beats ResNet50 by $\geq 5$ \Mpp{}
& Partially met. DenseNet121 wins on most CIFAR-10 L3 to L5 cells; TransNeXt-tiny wins on most MNIST cells; the per-cell margin is $\leq 5$ \Mpp{} in most cases, so no single architecture is unambiguously dominant. \\
\midrule
Inference latency
& $\leq 50$ ms per image on a standard GPU
& Met for ResNet50 and DenseNet121; close to the boundary for TransNeXt-tiny ($\sim 30$ ms with \texttt{torch.compile} disabled). \\
\midrule
Reproducibility
& Deterministic per-sample degradation
& Met. Byte-identity gate passes across seeds $\{42, 43, 44\}$. Multi-seed audit confirms $\sigma < 1.5$ \Mpp{} on every L3 headline cell. \\
\bottomrule
\end{tabular}
\end{table}

\begin{figure}[htbp]
  \centering
  \includegraphics[width=\textwidth]{\figpath{gantt/gantt_planned_vs_actual.png}}
  \caption{Planned (top) versus actual (bottom) project schedule. The actual schedule added three late-stage expansion phases (Phase D regularization sweep, Phase B2 / B2-nr / C2 THz-protocol decomposition, and the multi-seed audit) on top of the original plan.}
  \label{fig:gantt}
\end{figure}

\subsection{Success and failure analysis vs goals}
Of the four original goals from the preliminary specification:

\textbf{Goal 1 (high accuracy on degraded data, target 80\%).}\quad
\emph{Partially met.} The 80\% threshold is achieved on MNIST at every
level up to L3 and on CIFAR-10 at L1 and L2. At L3 on CIFAR-10 the best
cell ($\sim 74\%$) falls $\sim 6$ \Mpp{} short, and at L4 / L5 the
accuracy collapses on every backbone (the L5 ceiling on CIFAR-10 is
$\sim 40\%$). The shortfall is consistent with the per-axis attribution:
resolution alone costs $-27$ \Mpp{} at L3 to L5, which exceeds the
remaining headroom against the 80\% target.

\textbf{Goal 2 (minimize the accuracy drop, target $\leq 15$ \Mpp{}).}\quad
\emph{Partially met.} The drop is bounded by 15 \Mpp{} up to L2 on
CIFAR-10 and up to L3 on MNIST. Beyond those operating points the drop
grows monotonically with the degradation level, reaching $\sim 50$
\Mpp{} at L5 on CIFAR-10.

\textbf{Goal 3 (identify the optimal architecture).}\quad
\emph{Partially met.} A statistically clear $\geq 5$ \Mpp{} win for a
single architecture was not observed. DenseNet121 leads on most
CIFAR-10 L3 cells by 1 to 4 \Mpp{}; TransNeXt-tiny leads on most MNIST
cells by 1 to 3 \Mpp{}; the inter-architecture spread is comparable to
the multi-seed $\sigma$ in several cells. The honest reading of the
data is that the three architectures are within run-to-run noise of
each other at the operating points that matter, and that the choice of
architecture is secondary to the choice of degradation pipeline.

\textbf{Goal 4 (feasibility for edge implementation).}\quad
\emph{Met for the CNN backbones, marginal for TransNeXt-tiny.} Both
ResNet50 and DenseNet121 sit comfortably under the 50 ms per image
budget at \texttt{bf16-mixed} precision. TransNeXt-tiny is $\sim 30$ ms
per image with \texttt{torch.compile} disabled; enabling the graph
compiler (currently blocked by the foveal-attention path) would
likely halve the figure.

\subsection{Strengths of the chosen approach}
\begin{itemize}
  \item \textbf{Fair-comparison invariant.}\quad All three backbones see
        the same tensor shape ($224 \times 224$) and the same training
        recipe (modulo the per-paper LR convention), so the cross-backbone
        deltas are interpretable.
  \item \textbf{Deterministic degradation.}\quad The byte-identity gate
        across seeds $\{42, 43, 44\}$ means that the cross-phase deltas
        are not contaminated by data-side stochasticity.
  \item \textbf{Per-pipeline image-quality anchor.}\quad PSNR / SSIM
        per pipeline grounds the comparison in a physical scale; the
        Spearman rank correlations against accuracy (Figures
        \ref{fig:acc_vs_psnr}, \ref{fig:acc_vs_ssim}) are strong.
  \item \textbf{Multi-seed audit.}\quad The $\sigma < 1.5$ \Mpp{}
        bound on every L3 headline cell makes the cross-phase
        deltas interpretable.
\end{itemize}

\subsection{Weaknesses of the chosen approach}
\begin{itemize}
  \item \textbf{Synthetic degradation is not real THz.}\quad The
        five-axis pipeline is a mathematical model. Real THz physics
        includes wavelength-dependent diffraction, mixing-element
        nonlinearity, and detector-array dead pixels that the model
        does not capture. The conclusions transfer only to the extent
        that those secondary effects are dominated by the modeled
        axes.
  \item \textbf{Two small datasets.}\quad CIFAR-10 and MNIST are
        small-scale benchmarks. The generalization of the per-axis
        attribution to ImageNet-scale or to THz-domain object
        categories was not tested.
  \item \textbf{No real super-resolution baseline.}\quad The project
        evaluated training-time degradation exposure exclusively. A
        super-resolution baseline (a hypothetical Phase E) was scoped
        out at PDR for cost reasons. Without it the
        attribution-driven recommendation (invest in resolution
        restoration first) cannot be quantitatively compared
        against a real restoration pipeline.
  \item \textbf{No human-perception baseline.}\quad The acceptance
        criteria are model-side; a human upper bound (the percentage
        of L5 images a human reader can classify correctly) would
        contextualize the $\sim 40\%$ CIFAR-10 L5 ceiling.
\end{itemize}

\subsection{Recommendations for continuation}
\begin{enumerate}
  \item \textbf{Acquire real THz data.}\quad The synthetic pipeline is
        a proxy; a small real-world dataset (even 1000 images) at the
        target deployment frequency would let the next iteration
        calibrate the synthetic model against ground truth and
        renormalize the per-axis attribution.
  \item \textbf{Add a super-resolution pre-stage.}\quad Given that
        resolution dominates the per-axis attribution by $\sim 4 \times$
        (Phase C2), the next ablation should add a learned upsampler
        in front of each backbone and re-run Phase B2 / C2 with the
        new front end.
  \item \textbf{Enable graph compilation for TransNeXt.}\quad The
        foveal-attention path is currently not
        Inductor-graph-capturable; rewriting it as a compiled custom
        op would likely halve the inference latency and bring the
        TransNeXt backbone inside the 25 ms per image budget.
  \item \textbf{Robustness benchmark reproduction.}\quad Reproduce
        the per-axis attribution on the ImageNet-A robustness
        benchmark \cite{hendrycks2021naturaladv} to test whether the
        resolution-dominant finding holds at full ImageNet scale.
  \item \textbf{Calibration intervention.}\quad The ECE growth from
        $< 5$ \Mpp{} at L3 to $> 15$ \Mpp{} at L5 indicates that a
        post-hoc calibration step (temperature scaling) would
        improve the operational usefulness of the predictions
        without affecting top-1 accuracy.
\end{enumerate}

\subsection{Cost summary}
The actual operating cost of the campaign was dominated by GPU-hours
rather than by hardware or licensing. Table \ref{tab:cost} reports the
final budget. Operator-side hours are theoretical (NIS-equivalent of the
student-engineer market rate); they do not appear in the cash budget for
this project.

\begin{table}[htbp]
\centering
\caption{Final cost summary (NIS-equivalent, supervisor signature required for submission).}
\label{tab:cost}
\setlength{\tabcolsep}{4pt}
\renewcommand{\arraystretch}{1.1}
\begin{tabular}{lr}
\toprule
Item & Cost (NIS) \\
\midrule
Student engineer hours (2 students $\times$ 300 h $\times$ 60 NIS/h, theoretical) & 36\,000 \\
Supervisor hours (2 $\times$ 30 h $\times$ 250 NIS/h, theoretical) & 15\,000 \\
GPU compute (local RTX 5070, $\sim$ 56.5 + $\sim$ 48 + $\sim$ 24 GPU-hours)        & 0 (local) \\
Cloud GPU fallback budget (Colab Pro+, 8 months) & 400 \\
Office supplies and storage media & 500 \\
Subtotal & 51\,900 \\
Overhead (25\%) & 12\,975 \\
Contingency (10\%) & 5\,190 \\
\midrule
\textbf{Grand total} & \textbf{70\,065} \\
\bottomrule
\end{tabular}
\end{table}

\vspace{1em}
\noindent
\textit{Supervisor signature.\quad\dotfill\quad Date.\quad\dotfill}
